\documentclass[12pt]{article}

\usepackage[utf8]{inputenc}
\usepackage{amsmath, amssymb, amsthm}
\usepackage{geometry}
\geometry{a4paper, margin=2.5cm}

\title{Spații Produs și Măsuri Produs}
\author{}
\date{}

\theoremstyle{plain}
\newtheorem{theorem}{Teoremă}

\begin{document}

\maketitle

\section{Spații măsurabile produs}

Fie $(X, \mathcal{A})$ și $(Y, \mathcal{B})$ două spații măsurabile.
Definim \textbf{spațiul măsurabil produs} ca fiind:



\[
(X \times Y, \ \mathcal{A} \otimes \mathcal{B}),
\]



unde $\mathcal{A} \otimes \mathcal{B}$ este \textbf{sigma-algebra produs},
definită ca sigma-algebra generată de mulțimile dreptunghi:



\[
A \times B, \qquad A \in \mathcal{A}, \ B \in \mathcal{B}.
\]



Aceasta este cea mai mică sigma-algebră care conține toate dreptunghiurile
măsurabile.

\section{Măsura produs}

Fie $(X, \mathcal{A}, \mu)$ și $(Y, \mathcal{B}, \nu)$ două spații măsurabile
cu măsuri.

O măsură $\lambda$ pe spațiul produs


\[
(X \times Y, \ \mathcal{A} \otimes \mathcal{B})
\]


se numește \textbf{măsură produs} dacă pentru toate dreptunghiurile măsurabile:



\[
\lambda(A \times B) = \mu(A)\,\nu(B).
\]



\section{Teorema de existență și unicitate a măsurii produs}

\begin{theorem}[Teorema măsurii produs]
Fie $(X, \mathcal{A}, \mu)$ și $(Y, \mathcal{B}, \nu)$ două spații măsurabile
sigma-finite.

Atunci există o \textbf{unică} măsură $\mu \otimes \nu$ pe
$\mathcal{A} \otimes \mathcal{B}$ astfel încât:



\[
(\mu \otimes \nu)(A \times B) = \mu(A)\,\nu(B),
\qquad A \in \mathcal{A}, \ B \in \mathcal{B}.
\]



Această măsură se numește \textbf{măsura produs} a lui $\mu$ și $\nu$.
\end{theorem}

\section{Interpretare}

Teorema afirmă că:

\begin{itemize}
    \item dacă avem două spații măsurabile sigma-finite,
    \item atunci putem construi în mod riguros o măsură pe spațiul lor produs,
    \item iar această măsură este unică și se comportă corect pe dreptunghiuri.
\end{itemize}

Aceasta este baza pentru integralele multiple, teorema lui Fubini și
construcția distribuțiilor comune ale variabilelor aleatoare.

\section{Exemplu}

Pe $\mathbb{R}^2$, dacă $\mu = \nu = \lambda$ este măsura Lebesgue pe
$\mathbb{R}$, atunci măsura produs:



\[
\lambda \otimes \lambda
\]



este măsura Lebesgue bidimensională (măsura ariei).

\end{document}
